%!TeX encoding = UTF-8 Unicode
%!TeX spellcheck = en-UK
% !TeX root = Perrinet26mnesis.tex
% -*- coding: UTF-8; -*-
%%%-----------------------------------------------------------------
\begin{figure}[!ht]
\begin{center}
\includegraphics[width=\textwidth]{retrieval.pdf}
\end{center}
\caption{Sequential retrieval of the $M=16$ stored motifs. The motifs are concatenated into a single stream, each flanked on both sides by $200$ silent steps, and only their first $D=41$ steps are given to the network as a trigger. The $F_1$ score (green), computed in consecutive $60$-step windows against the frozen motifs, is shown over the duration of each motif (it is not shown in the silent intervals, where the target is empty), together with the trigger mask, equal to $1$ on each trigger window (orange), and the ideal retrieval envelope, equal to $1$ over the full $T=1000$-step duration of each motif (blue). The motifs are numbered above their envelope, and the number of the second one, shown in Figure~\ref{fig:target}, is in bold. Retrieval succeeds for $12$ of the $16$ motifs: the $F_1$ score rises at the trigger onset and stays close to $1$ until the end of the motif, in two cases (2nd and 11th motifs) after a partial dip that follows the trigger. For the 4th, 10th, 12th and 13th motifs, the replay starts correctly but breaks down after about $200$ to $600$ steps: the $F_1$ score then collapses to $0$ and does not recover before the end of the motif. Grey dashed lines indicate extremal $F_1$ values $0$ or $1$.}\label{fig:retrieval}
\end{figure}
%%%-----------------------------------------------------------------
